\section{Research Questions}
\begin{description}
\item[RQ1] How often do coding agents implement requested backend changes while preserving required invariants?
\item[RQ2] Does explicit invariant-oriented guidance improve joint correctness relative to standard instructions?
\item[RQ3] Which categories of invariant violation occur most frequently?
\item[RQ4, exploratory] How do execution time, tokens and tool interactions relate to correctness?
\end{description}
The preliminary hypothesis predicts improved correctness with invariant-oriented guidance. It is
untested. The primary planned estimand is the difference in success probability between conditions,
equally weighted across validated tasks and selected systems.

\section{Methodology}
The draft design crosses up to eight validated backend modification tasks, two agent systems, two
prompt conditions and three independent executions per cell. The ceiling of 96 main executions is
provisional and is not a statistically justified sample size. A separate pilot will reassess task count,
repetitions, feasibility, power and detectable effects before main collection.

Both conditions receive identical feature requests, pinned source, documentation access and underlying
requirements. Condition B adds structured attention to affected state transitions, transaction
boundaries, retries and verification. Prompt length and planning attention are part of this bundled
intervention. Independent sessions and clean environments will prevent cross-condition carryover.

Each accepted task will require a reproducible baseline, clear new behavior, a human-reviewed reference,
automated feature and invariant checks, and a relevant deliberately defective implementation rejected
by the tests. Hidden evaluations remain outside the agent's accessible filesystem. The primary outcome is
\[
\texttt{full\_task\_success} =
\texttt{feature\_passes} \land \texttt{all\_required\_invariants\_pass}.
\]
Individual outcomes and histories will be retained. Environment failures are investigated separately
and do not silently become ordinary correctness failures. Missing measurements remain missing.

Analyses will describe task/system/condition rates and matched condition differences. Repeated runs
are nested within tasks; marginal intervals cannot establish cluster-aware uncertainty. Mixed-effects
logistic regression will be considered only if supported by adequate variation and identifiable data.
Sparse data may justify descriptive conclusions alone. Resource associations remain exploratory.

\section{Experimental Setup}
The intended subject is Dragon Market\footnote{\url{https://github.com/houtanrocky/dragon-market}}.
The inspected candidate baseline is commit
\texttt{162c5cab9fa886445dd102f4a99961027f90aebb}. Its Go build, test, race-enabled test and vet
commands passed using Go 1.26.0 on Windows with isolated Testcontainers PostgreSQL. A targeted
characterization test also reproduced a pre-existing daily-spending defect involving release of a
prior-day reservation. Its passing outcome does not establish preservation of the intended daily bound;
that requirement needs an explicit scope or reviewed baseline decision before main collection.
These are setup checks, not coding-agent observations.

Candidate tasks cover refunds, settlement, purchases, transfers, reservations, rollback, limits and
audit consistency. None is currently validated. Inspection found existing settlement, purchase-safety,
cancellation, rollback and audit implementations, so those candidates need meaningful feature extensions.
The baseline also contains public research-readiness material and a workbench; the exact content exposed
to agents must be frozen identically across conditions and separated from new hidden study oracles.

\textbf{Setup fields to complete after pilot:} exact subject/research commits, accepted task inventory,
agent and model versions, prompt hashes, budgets/retries, tool permissions, OS/Go/dependency versions,
PostgreSQL image digest and isolation configuration, seed/reset/schedule policy, baseline failures,
and evaluator commands. No coding-agent integration, validated benchmark or empirical outcome is implied.
